\chapter{No es bueno estar solo}

\epigraph{«Varón y mujer los creó.»}{Génesis 1,27}

Hay un texto que casi todo el mundo cree conocer, y que casi nadie ha leído con cuidado.

Me refiero, claro, al relato de la creación. Todos tenemos en la cabeza una especie de resumen que va más o menos así: Dios crea al hombre, lo pone en el jardín, ve que está solo, le quita una costilla y fabrica a la mujer. Fin. La mujer aparece después, como un añadido, casi como un parche. Como si Dios hubiera hecho al varón, se hubiera dado cuenta de que le faltaba algo, y hubiera improvisado una solución.

Pero eso no es lo que dice el Génesis. O, mejor dicho, no es \emph{todo} lo que dice. Porque el Génesis tiene dos relatos de la creación, no uno. Y cuando los leemos juntos, con calma, lo que encontramos es bastante más hondo ---y bastante más hermoso--- de lo que el resumen sugiere.

\section{Lo que el Génesis dice de verdad}

El primer relato está en el capítulo 1. Es el más solemne, el más cósmico. Dios va creando por etapas ---la luz, el firmamento, las aguas, los animales--- y cada vez dice: «Es bueno.» Hasta que llega al ser humano. Y entonces el texto hace algo distinto. No dice «hágase el hombre» como dijo «hágase la luz». Dice algo mucho más íntimo:

\begin{quotation}
\textit{«Hagamos al hombre a nuestra imagen y semejanza. [\ldots] Y creó Dios al hombre a su imagen, a imagen de Dios lo creó; varón y mujer los creó.»}\footnote{Gn 1,26-27. El término hebreo \textit{adam} no designa aquí al varón, sino al ser humano, a la humanidad. Solo en el segundo relato (Gn 2) se introduce la diferenciación \textit{ish} (varón) e \textit{ishshah} (mujer).}
\end{quotation}

Quiero que nos detengamos aquí un momento, porque esta frase contiene algo que cambia toda la lectura. Fijémonos: Dios no crea primero al varón y después a la mujer. Crea al ser humano ---\textit{adam}, la humanidad--- y lo crea «varón y mujer». Ambos a la vez. Ambos imagen de Dios. No hay un primero y un segundo. No hay un original y una copia. Hay un ser humano que, desde el principio, es relación: varón \emph{y} mujer.

Solo después de crear así al ser humano, Dios dice algo que no había dicho de nada en toda la creación: «No es bueno que el hombre esté solo.»\footnote{Gn 2,18. Es la primera vez en todo el relato de la creación que algo «no es bueno». Juan Pablo II desarrolla extensamente el significado de esta «soledad originaria» en sus catequesis sobre la Teología del Cuerpo. Cf. Juan Pablo II, \textit{Hombre y mujer los creó. Catequesis sobre el amor humano}, audiencias generales de septiembre a noviembre de 1979.}

Y aquí está la clave: si la mujer fuera simplemente un complemento que viene a llenar un hueco, a solucionar un problema de diseño, Dios habría creado primero al varón solo y luego habría tenido que corregir su propia obra. Pero no es eso lo que pasa. Lo que pasa es que el ser humano es, desde su misma creación, un ser relacional. La capacidad de amar, la necesidad del otro, no es algo que se le añade después: está inscrita en su origen. Es constitutiva. La mujer no viene a arreglar un defecto del varón. Ni el varón a completar una carencia de la mujer. El «no es bueno que el hombre esté solo» no es la constatación de un fallo, sino la revelación de una verdad: estamos hechos para el encuentro.

\section{La biología ya lo sabía}

Si habéis leído la primera parte de este libro, todo esto os sonará. Y es que resulta que la fe y la ciencia, cada una desde su lenguaje, dicen lo mismo.

Hemos visto que nacemos radicalmente inmaduros: un «feto externo» que necesita meses de cuidado constante para sobrevivir.\footnote{Cf. capítulo 1 de este libro.} Hemos visto que nuestro cuerpo está equipado con mecanismos hormonales ---oxitocina, vasopresina, dopamina--- que nos empujan a vincularnos, a cuidar, a quedarnos.\footnote{Cf. capítulos 2 y 3.} Hemos visto que la diferencia entre varón y mujer no es un accidente cultural, sino un dato inscrito en nuestra biología que hace posible la complementariedad y el encuentro.\footnote{Cf. capítulo 5.}

La biología nos dice: no estáis hechos para vivir solos. Necesitáis al otro. El amor no es un lujo, es una condición de supervivencia.

Y ahora el Génesis llega y le pone nombre a lo que el cuerpo ya sabía: «No es bueno que el hombre esté solo.» No es bueno, no porque le falte compañía, sino porque sin el otro no puede amar. Y sin amar, no puede ser plenamente lo que es: imagen de Dios.

\section{Imagen de un Dios que es comunión}

Aquí entramos en algo que la tradición cristiana ha ido comprendiendo poco a poco a lo largo de siglos, y que quizá sea una de las intuiciones más profundas de toda la teología.

Cuando el Génesis dice que el ser humano es «imagen de Dios», ¿qué quiere decir? ¿Que nos parecemos a Dios físicamente? Obviamente no. ¿Que somos inteligentes, libres, creativos? Sí, también. Pero hay algo más radical.

Dios, para la fe cristiana, no es un ser solitario. Dios es Trinidad: Padre, Hijo y Espíritu Santo. No tres dioses, sino un solo Dios que es, en sí mismo, relación, comunión, entrega. El Padre se entrega totalmente al Hijo. El Hijo se entrega totalmente al Padre. Y esa entrega mutua es tan real, tan fecunda, que «produce» una tercera persona: el Espíritu Santo, que es el Amor en persona.\footnote{«Dios, que es amor y vive en sí mismo un misterio de comunión personal de amor, crea al género humano a su imagen [\ldots] inscribiendo en la humanidad del hombre y de la mujer la vocación y, consiguientemente, la capacidad y la responsabilidad del amor y de la comunión.» Juan Pablo II, Exhortación apostólica \textit{Familiaris Consortio}, n. 11 (1981).}

Dios no \emph{tiene} amor. Dios \emph{es} amor.\footnote{1~Jn 4,8.} Y ser imagen de ese Dios significa estar hecho para amar como Él ama: entregándose, saliendo de uno mismo, haciéndose don para el otro.

Por eso el ser humano no puede realizarse en solitario. No es una cuestión de debilidad o de dependencia psicológica. Es una cuestión de identidad. Somos imagen de un Dios-comunión: solo amando nos parecemos a Él. Solo en la entrega al otro encontramos quiénes somos.

\section{La costilla, la carne, el grito}

Ahora sí podemos entender el segundo relato del Génesis ---el de la costilla--- en toda su profundidad. Porque no es un cuento ingenuo sobre cirugía divina. Es una de las páginas más densas jamás escritas sobre lo que significa encontrar al otro.

Dios hace desfilar ante el hombre a todos los animales para que les ponga nombre.\footnote{Gn 2,19-20.} En la mentalidad bíblica, poner nombre es conocer, es comprender la esencia de algo. El hombre nombra a cada animal, los conoce, los entiende. Pero en ninguno se reconoce. Ninguno es «como él». Hay una soledad que ningún animal puede llenar ---por mucho que queramos a nuestro perro, que conste---.\footnote{Juan Pablo II llama a esta experiencia la «soledad originaria»: el ser humano se descubre distinto de todo lo demás, se reconoce como persona, como ser que busca un «tú». Cf. Juan Pablo II, audiencia general del 10 de octubre de 1979.}

Y entonces Dios actúa. No toma barro del suelo, como hizo para formar al primer ser humano. Toma algo del hombre mismo: una costilla, que en hebreo (\textit{tsela}) puede significar también «costado», «lado».\footnote{El término hebreo \textit{tsela'} aparece en otros libros del Antiguo Testamento con el significado de «lado» (por ejemplo, en la descripción de los laterales del Arca de la Alianza, Ex 25,12). La traducción como «costilla» es posible, pero empobrece el sentido: lo que Dios toma no es un hueso, sino un lado, una mitad.} La imagen es poderosa: la mujer no viene de fuera, no es un cuerpo extraño. Sale del mismo ser del hombre. Comparte su sustancia, su dignidad, su naturaleza.

Y cuando el hombre despierta y la ve, se produce el primer poema de la historia humana. Las primeras palabras que un ser humano pronuncia en toda la Biblia no son una oración, ni un mandamiento, ni una queja. Son un grito de asombro, de reconocimiento, de alegría:

\begin{quotation}
\textit{«¡Esta sí es hueso de mis huesos y carne de mi carne!»}\footnote{Gn 2,23.}
\end{quotation}

\textit{Esta sí.} En hebreo, \textit{zo't happa'am} ---«¡por fin!», «¡esta vez sí!»---. Después de haber visto pasar a todos los animales sin reconocerse en ninguno, el hombre ve a la mujer y por fin se encuentra. No encuentra a alguien igual que él ---si fuera igual, no habría encuentro---. Encuentra a alguien que es distinto y a la vez de su misma carne. Alguien a quien puede mirar cara a cara. Alguien a quien puede amar.

Hay un detalle que pasa desapercibido y que es precioso: el hombre solo se reconoce a sí mismo como varón (\textit{ish}) cuando reconoce a la mujer (\textit{ishshah}). Antes de ella, no tenía nombre propio. Era simplemente \textit{adam}, el ser humano. Solo ante el otro descubre quién es. Solo en la relación se revela la identidad.\footnote{«El hombre sólo se define como ``varón'' ante la ``mujer''; sólo la presencia del ``tú'' revela plenamente su ``yo''.» Cf. Juan Pablo II, audiencia general del 7 de noviembre de 1979.}

\section{El matrimonio como sacramento}

«Por eso dejará el hombre a su padre y a su madre, se unirá a su mujer, y serán una sola carne.»\footnote{Gn 2,24. Jesús mismo citará este versículo en Mt 19,5 para fundamentar la indisolubilidad del matrimonio, remitiendo a sus oyentes no a la Ley de Moisés, sino «al principio».}

Así termina el relato del Génesis. Y así comienza la teología del matrimonio.

Cuando dos personas se casan por la Iglesia, hay algo que mucha gente no sabe, o que sabe pero no ha terminado de entender: los esposos no «reciben» el sacramento del matrimonio como quien recibe un paquete. Los esposos \emph{son} los ministros del sacramento.\footnote{«Los esposos, como ministros de la gracia de Cristo, se confieren mutuamente el sacramento del Matrimonio, expresando ante la Iglesia su consentimiento.» \textit{Catecismo de la Iglesia Católica}, n. 1623. A diferencia de los demás sacramentos, el ministro del matrimonio no es el sacerdote o el diácono, sino los propios contrayentes. El sacerdote es testigo cualificado de la Iglesia, pero no administra el sacramento.} Son ellos quienes se lo dan mutuamente. El sacerdote no casa a nadie: es testigo. Testigo cualificado, sí, en nombre de la Iglesia. Pero quien realiza el sacramento son el hombre y la mujer que se dicen el uno al otro: «Yo te recibo como esposo. Yo te recibo como esposa.»

Esto no es un detalle litúrgico. Es una revolución.

Significa que el matrimonio no es un trámite administrativo con flores y arras. No es una bendición decorativa que la Iglesia añade a lo que la pareja ya tenía. Es un acto en el que dos personas, con su libertad y ante Dios, se convierten en signo vivo de algo mucho más grande que ellas mismas: el amor de Cristo por su Iglesia.\footnote{Ef 5,31-32: «"Por eso dejará el hombre a su padre y a su madre y se unirá a su mujer, y los dos serán una sola carne." Gran misterio es este, y yo lo refiero a Cristo y a la Iglesia.»}

Y ese signo no es abstracto. Ese signo es la vida entera de los esposos. Cada gesto de cuidado, cada perdón concedido, cada noche en vela junto a un hijo enfermo, cada renuncia silenciosa, cada abrazo después de una discusión: todo eso es sacramento. Todo eso hace visible, aquí y ahora, un amor que viene de más arriba.

\subsection*{Alianza, no contrato}

Hay una diferencia fundamental entre un contrato y una alianza, y entenderla cambia la forma de mirar el matrimonio.\footnote{Cf. Scott Hahn, \textit{La cena del Cordero}, Rialp, 2001. Hahn desarrolla con claridad la diferencia entre contrato y alianza en el contexto bíblico.}

Un contrato intercambia bienes y servicios: yo te doy esto, tú me das aquello. Si una parte incumple, la otra queda libre. Un contrato se puede rescindir. Un contrato protege intereses.

Una alianza intercambia personas: yo me doy a ti, tú te das a mí. No hay letra pequeña. No hay cláusula de rescisión. Una alianza crea un vínculo que transforma a quienes la contraen. «Ya no son dos, sino una sola carne»,\footnote{Mt 19,6.} dice Jesús. No dijo «un solo patrimonio» ni «un solo proyecto». Dijo «una sola carne»: una sola vida entregada.

El matrimonio cristiano es una alianza. Y cuando los novios se dicen «sí», están haciendo lo mismo ---salvando las infinitas distancias--- que hizo Dios con su pueblo: vincularse para siempre, sin condiciones, por amor. «Yo seré tu Dios y tú serás mi pueblo.»\footnote{Jr 31,33; cf. Os 2,21-22: «Te desposaré conmigo para siempre; te desposaré en justicia y en derecho, en amor y en compasión.»} Los profetas, de hecho, usaron constantemente la imagen del matrimonio para hablar de la relación entre Dios e Israel.\footnote{Toda la tradición profética emplea la metáfora nupcial: Oseas, Isaías, Jeremías, Ezequiel. El Cantar de los Cantares se ha leído, desde antiguo, como canto del amor entre Dios y su pueblo. No es casualidad que Jesús inaugure su vida pública en una boda (Jn 2,1-11).} No al revés. No es que el matrimonio sea «como» la alianza de Dios. Es que la alianza de Dios se expresa como un matrimonio, porque el matrimonio es la imagen más cercana que tenemos los humanos del amor total.

\section{Una cruz a la que abrazar}

Quiero terminar este capítulo con algo que no está en ningún manual de teología, pero que a mí me ha dado mucho que pensar. Es algo que suelo compartir en las charlas con novios, y que siempre genera un silencio especial en la sala.

Jesús dijo: «Nadie va al Padre sino por mí.»\footnote{Jn 14,6.} Y también dijo: «El que quiera seguirme, que tome su cruz y me siga.»\footnote{Cf. Mt 16,24; Mc 8,34; Lc 9,23.}

Normalmente, cuando pensamos en la Cruz, pensamos en sufrimiento. En sacrificio. En algo pesado, ominoso, doloroso. «Llevar la cruz» se ha convertido en sinónimo de aguantar, de soportar, de cargar con algo que preferirías no tener.

Pero paremos un momento. ¿Qué es la Cruz, en realidad?

La Cruz es el lugar donde Cristo entrega su vida. Sí, con dolor. Sí, con sufrimiento. Pero la Cruz no es principalmente sufrimiento. La Cruz es principalmente \emph{entrega}. Es el lugar donde Jesús dice, con todo su cuerpo y toda su sangre: «Te amo hasta el final. Doy mi vida por ti.»\footnote{Cf. Jn 13,1: «Habiendo amado a los suyos que estaban en el mundo, los amó hasta el extremo.»}

Y ahora pensemos en el matrimonio.

¿No es el matrimonio, precisamente, el lugar donde dos personas entregan su vida por amor? No de golpe, no con un gesto dramático, sino día a día, gota a gota. En el café que preparas por la mañana sabiendo cómo le gusta. En el cansancio compartido después de acostar a los niños. En la renuncia a tener razón para poder tener paz. En el «me quedo contigo» de cada día, que es menos espectacular que el del día de la boda pero infinitamente más real.

Si la Cruz es entrega de vida por amor, y el matrimonio es entrega de vida por amor\ldots{} entonces mi mujer es mi cruz. Mi marido es mi cruz. No en el sentido de «qué cruz de marido me ha tocado» ---que también tiene su gracia---. En el sentido más profundo: la persona a la que amo es la razón por la que entrego mi vida. Es mi camino de santidad. Es mi puerta al Padre.

«El que quiera seguirme, que tome su cruz.» Y Dios, que nos conoce bien, nos lo pone fácil. Nos da una cruz a la que podemos abrazar cada noche. Una cruz que nos devuelve el abrazo. Una cruz que nos dice «te quiero» y que nos necesita tanto como nosotros a ella. Y aquí está la clave: que devuelve el abrazo. Una cruz que aplasta, que humilla, que hace daño, no es la cruz de Cristo: es la cruz del verdugo. La cruz del matrimonio es siempre entrega mutua, nunca sufrimiento impuesto por uno sobre el otro.\footnote{Sobre las situaciones de violencia o abuso en el matrimonio, véase el capítulo 8 de este libro.}

Hay gente que busca la santidad en grandes penitencias, en ayunos heroicos, en gestas extraordinarias. Y está bien. Pero para los casados, el camino es más sencillo y más difícil a la vez: amar a la persona que tienes al lado. Amarla cuando es fácil y cuando no lo es. Amarla cuando te enamora y cuando te saca de quicio. Amarla como Cristo ama a su Iglesia: «entregándose por ella».\footnote{Ef 5,25.}

Esa es la cruz del matrimonio. Y es, quizá, la forma más hermosa de llevar la cruz que existe.

\paracaminar

\begin{enumerate}
  \item Leed juntos Génesis 1,26-27 y Génesis 2,18-24. ¿Qué os dice sobre vuestro matrimonio que el ser humano sea, desde el principio, «varón y mujer», creado para la relación?

  \item ¿Habéis pensado alguna vez en vuestro matrimonio como una \emph{alianza} y no como un contrato? ¿Qué diferencia práctica tiene eso en vuestra vida diaria?

  \item La imagen de la Cruz como entrega de amor: ¿cómo la recibís? ¿Podéis identificar momentos concretos en los que «entregar la vida» por el otro os haya hecho crecer como personas?
\end{enumerate}

\vinetacapitulo{ch07}
